\section*{Hoofdstuk \ref{ch:b1:electronic-effects} --- Elektronische effecten en reactieve intermediairen}

\begin{solution}{exo:b1:electronic-effects:1}
2-Methylbutaan, \ce{CH3-CH(CH3)-CH2-CH3}: de drie \ce{CH3} zijn primair,
de \ce{CH2} secundair, de CH tertiair. 2-Broom-2-methylpropaan: de drie
\ce{CH3} zijn primair, het centrale koolstofatoom tertiair. Van de vier
bromiden \ce{C4H9Br} (1-broombutaan, 2-broombutaan,
1-broom-2-methylpropaan, 2-broom-2-methylpropaan) geeft het laatste het
tertiaire \omterm{def:b1:electronic-effects:intermediates}{carbokation}, het stabielste.
\end{solution}

\begin{solution}{exo:b1:electronic-effects:2}
$10^{4.76 - 2.87} = 10^{1.89} = 78$; $10^{4.76 - 0.52} = 10^{4.24} =
\num{1.7e4}$.
\end{solution}


\begin{solution}{exo:b1:electronic-effects:3}
Ethanoaat: de C=O en de \ce{C-O^-} wisselen tussen de twee
zuurstofatomen. 4-Nitrofenoxide: het \omterm{def:b1:lewis-resonance-vsepr:lewis}{vrije paar} van zuurstof vormt een
C=O, de dubbele bindingen van de ring verschuiven, het koolstofatoom dat
\ce{NO2} draagt vormt een C=N-binding, en een N=O-paar gaat naar zuurstof:
een structuur met C=O aan het ene uiteinde van de ring, C=N aan het
andere en beide zuurstofatomen van de nitrogroep negatief (een chinoïde
structuur). Allylkation:
\ce{CH2=CH-CH2+} $\leftrightarrow$ \ce{+CH2-CH=CH2}.
\end{solution}

\begin{solution}{exo:b1:electronic-effects:4}
Tweede pijl: van de O--H-binding van water naar zijn zuurstofatoom;
producten \ce{H2O} (uit \ce{OH-}, nu neutraal) en \ce{OH-} (de zuurstof
van het vroegere water houdt het paar, lading $-1$). Voor ammoniak: een
pijl van het \omterm{def:b1:lewis-resonance-vsepr:lewis}{vrije paar} van N naar een H van \ce{H3O+}, en een van die
O--H-binding naar de zuurstof: N wordt $+1$, O wordt 0.
\end{solution}

\begin{solution}{exo:b1:electronic-effects:5}
Ethanol (15.9) $<$ water (14.0) $<$ fenol (9.99, delokalisatie in de
ring) $<$ 4-nitrofenol (7.16, \ce{NO2} neemt de lading op) $<$
ethaanzuur (4.76, twee gelijkwaardige zuurstofatomen) $<$
trichloorethaanzuur (0.51, drie $-I$-chlooratomen).
\end{solution}

\begin{solution}{exo:b1:electronic-effects:6}
Aniline (4.60) $<$ pyridine (5.23) $<$ ammoniak (9.25) $<$ methylamine
(10.66). In aniline wordt het \omterm{def:b1:lewis-resonance-vsepr:lewis}{vrije paar} met de ring gedeeld (structuren
met C=N$^+$ en de negatieve lading op \textit{ortho} en
\textit{para}); N protoneren zou die delokalisatie kosten.
\end{solution}

\begin{solution}{exo:b1:electronic-effects:7}
In 3-nitrofenoxide leggen de \omterm{def:b1:lewis-resonance-vsepr:resonance}{resonantiestructuren} de lading op de
koolstofatomen \textit{ortho} en \textit{para} ten opzichte van de
zuurstof, nooit op het koolstofatoom dat \ce{NO2} draagt: geen mesomere
hulp, dus zwakker dan het 4-isomeer. De nitrogroep trekt nog steeds
elektronen via de $\sigma$-bindingen ($-I$), dus sterker dan fenol.
\end{solution}

\begin{solution}{exo:b1:electronic-effects:8}
\ce{CH3+} $<$ \ce{CH3CH2+} $<$ \ce{CH2=CH-CH2+} $<$ \ce{C6H5CH2+}
$\approx$ \ce{(CH3)3C+} $<$ \ce{CH3OCH2+}: alkylgroepen geven elektronen;
allyl en benzyl delokaliseren de lading (twee en vier structuren); het
\omterm{def:b1:lewis-resonance-vsepr:lewis}{vrije paar} van zuurstof geeft een structuur met overal een volledig
octet. De volgorde in het midden van de lijst hangt af van het milieu.
\end{solution}

\begin{solution}{exo:b1:electronic-effects:9}
\ce{OH-} (\omterm{def:b1:electronic-effects:nucleophile}{nucleofiel}) valt het koolstofatoom van \ce{CH3Br} (\omterm{def:b1:electronic-effects:nucleophile}{elektrofiel})
aan, waarbij het C--Br-paar op Br vertrekt. \ce{H2O} geeft een paar aan
\ce{H+} (hier spreekt men van zuur en base). \ce{CN-} valt het
carbonylkoolstofatoom van ethanal aan, waarbij het $\pi$-paar van C=O naar
zuurstof gaat. \ce{NH3} geeft zijn \omterm{def:b1:lewis-resonance-vsepr:lewis}{vrije paar} aan de lege orbitaal van
boor in \ce{BF3}.
\end{solution}

\begin{solution}{exo:b1:electronic-effects:10}
Nee: $K_a = 10^{-0.51} = 0.31$ is veel groter dan $C$; het zuur is bijna
volledig gedissocieerd. $x^2/(0.010 - x) = 0.31$ geeft
$x = \qty{9.7e-3}{mol/L}$, $\mathrm{pH} = 2.01$, \omterm{def:b1:predominant-reaction:dissociation}{dissociatiegraad}
\qty{97}{\percent}.
\end{solution}

\begin{solution}{exo:b1:electronic-effects:11}
Alkylgroepen stuwen elektronen ($+I$) naar het zuurstofatoom van het
alkoxide en maken zijn lading minder stabiel, des te meer naarmate er
meer groepen zijn. Grotere alkoxiden worden ook minder goed door water
gesolvateerd, dat een klein geladen zuurstofatoom beter stabiliseert dan
een dat tussen alkylgroepen begraven ligt. Beide verhogen de
$\mathrm{p}K_a$.
\end{solution}

\begin{solution}{exo:b1:electronic-effects:12}
In een amide wordt het \omterm{def:b1:lewis-resonance-vsepr:lewis}{vrije paar} van stikstof gedeeld met de C=O
(structuur \ce{N+=C-O^-}); het gebruiken voor een proton of een
\omterm{def:b1:electronic-effects:nucleophile}{elektrofiel} zou die stabilisatie tenietdoen. In \omterm{def:b1:predominant-reaction:strong-acid}{sterk zuur} wordt het amide
geprotoneerd op de zuurstof, die de negatieve partiële lading draagt.
\end{solution}

\begin{solution}{pb:b1:electronic-effects:1}
\textbf{1.} $10^{1.89} = 78$; $10^{1.52} = 33$; $10^{0.84} = 6.9$.
\textbf{2.} Nee: elk volgend chlooratoom helpt minder. De lading van het
carboxylaat is al goed gespreid, en de metingen worden minder nauwkeurig
naarmate het zuur de volledige dissociatie nadert.
\textbf{3.} Elk Cl trekt elektronendichtheid aan via de C--C- en
C--O-bindingen en stabiliseert het carboxylaat ($-I$).
\textbf{4.} Het \omterm{def:b1:electronic-effects:inductive}{inductieve effect} dooft uit binnen twee of drie
bindingen.
\textbf{5.} De twee zuren lijken even sterk (0.52 en 0.51). In water zijn
beide bij gebruikelijke concentraties bijna volledig gedissocieerd
(\cref{ch:b1:predominant-reaction}, nivellering): een $\mathrm{p}K_a$ rond
0 is de grens van wat water kan onderscheiden, en de twee waarden hebben
onzekerheden van enkele tienden.
\textbf{6.} Ethaanzuur: zure vorm (\qty{98}{\percent}); chloorethaanzuur:
beide, met het anion licht overheersend ($10^{3 - 2.87} = 1.3$); de drie
andere: anion.
\textbf{7.} \ce{CH3-C(=O)-O^-} $\leftrightarrow$ \ce{CH3-C(-O^-)=O}.
\textbf{8.} Twee gelijke C--O-bindingen, tussen een dubbele en een enkele
binding in.
\textbf{9.} De lading zit op het enige zuurstofatoom, zonder
$\pi$-binding om ze te delen.
\textbf{10.} $10^{15.9 - 4.76} = 10^{11.1}$, ongeveer $10^{11}$.
\textbf{11.} De lading van het fenoxide spreidt zich naar drie
ringkoolstofatomen, die minder elektronegatief zijn dan zuurstof: een
winst, kleiner dan in een carboxylaat: $\mathrm{p}K_a$ 9.99 tussen 4.76
en 15.9.
\textbf{12.} $K = 10^{6.37 - \mathrm{p}K_a}$: ethaanzuur $10^{1.61} = 41$,
begunstigd (er komt koolstofdioxide vrij met een overmaat
waterstofcarbonaat); fenol $10^{-3.62}$ en ethanol $10^{-9.5}$: geen
reactie.
\textbf{13.} 4-Nitrofenol (7.16) $>$ 2-nitrofenol (7.23) $>$ 3-nitrofenol
(8.36) $>$ fenol (9.99).
\textbf{14.} De chinoïde structuur van \cref{exo:b1:electronic-effects:3},
met de lading op een zuurstofatoom van de nitrogroep.
\textbf{15.} De lading bereikt alleen de koolstofatomen \textit{ortho} en
\textit{para} ten opzichte van de zuurstof; het nitrokoolstofatoom staat
\textit{meta}.
\textbf{16.} Het $-I$-effect van \ce{NO2}.
\textbf{17.} In 2-nitrofenol vormt de OH een \omterm{def:b1:intermolecular-forces-solvents:hbond}{waterstofbrug} met een
zuurstofatoom van de naburige nitrogroep, wat de zure vorm stabiliseert
en het proton iets moeilijker verwijderbaar maakt.
\textbf{18.} Anionfracties $1/(1 + 10^{\mathrm{p}K_a - 7.5})$: 4-nitro
\qty{69}{\percent}, 2-nitro \qty{65}{\percent}, 3-nitro
\qty{12}{\percent}. Het zuur en het anion van 4-nitrofenol absorberen
verschillend, het anion in het zichtbare: zijn kleur verschijnt rond zijn
$\mathrm{p}K_a$.
\textbf{19.} $x^2/(0.010 - x) = 10^{-0.52} = 0.30$: $x = \qty{9.7e-3}{mol/L}$,
$\mathrm{pH} = 2.01$: bijna een \omterm{def:b1:predominant-reaction:strong-acid}{sterk zuur}.
\textbf{20.} $\frac12(4.76 + 2.00) = 3.38$.
\textbf{21.} \ce{CF3COOH + CH3COO- -> CF3COO- + CH3COOH}, $K = 10^{4.76 -
0.52} = 10^{4.24}$: \omterm{def:b1:extent-q-and-k:quantitative}{kwantitatief}.
\textbf{22.} $K_a(\ce{CF3COOH})/K_a(\ce{CH3COOH}) =$ \textbf{$10^{4.24} =
\num{1.7e4}$}, ongeveer tienduizend.
\end{solution}
